\documentclass[11pt]{article}
\usepackage[margin=1.1in]{geometry}
\usepackage{amsmath,amssymb,amsthm}
\usepackage{hyperref}
\usepackage{enumitem}

\newtheorem{theorem}{Theorem}
\newtheorem{lemma}{Lemma}
\newtheorem{corollary}{Corollary}
\newtheorem{proposition}{Proposition}
\newtheorem{remark}{Remark}
\newtheorem{definition}{Definition}

\title{Cycle 2, Part I: The Weil-Forcing Frame --- Geometric Zero-or-Full Dichotomy Completed,\\
Input Stratification (Monoid Layer, Torsion Gap),\\
and the End-Mass Law of the Exact Carrier}
\author{GST Hodge Campaign (Mode B, cycle 2)}
\date{25 October 2026}

\begin{document}
\maketitle

\noindent\textbf{Base checkpoint:} branch \texttt{research/gpt6-native-point-primitive-span-20261017} at
$2ac3b2f$. This derivation continues
\texttt{2026-10-24-gst-forcing-site-pure-realization-congruence-closure-and-gr25-calibration} (Theorem A, Lemma B1,
Theorem B2) and \texttt{2026-10-24-gst-gr25-lattice-completion-pure-pair-and-universal-two-center-descent-law}
(the two-center descent law, all $p$), both adversarially audited to PASS (2026-10-25), and the Weil-fourfold records
of 2026-10-17/18/19/21. It is part of campaign cycle 2 (the forcing input itself). External named theorems are used only
with hypotheses stated; every external leaf is declared in \S\ref{sec:scope}. No Lean file is changed by this document;
the companion Lean certificate \texttt{GSTWeilForcingFrameCertificate.lean} certifies the pure-algebra cores of
\S\ref{sec:correspondence} and \S\ref{sec:endsignature} on GitHub Actions.

\begin{center}
\fbox{\parbox{0.93\textwidth}{
\textbf{What this document proves.}
(1) The geometric zero-or-full dichotomy of the Weil-fourfold discriminator is completed to a full frame: the
Weil projector is exhibited as an \emph{explicit rational algebraic correspondence} with verified coefficients, its
invariance and image facts are reduced to kernel-verified algebra, and the two outcomes are characterized exactly:
$S=0$ forces every Weil class nonalgebraic and forbids \emph{every} holomorphic realization of the exact Weil carrier
at \emph{every} scale; $S=W_\mathbb{Q}$ produces \emph{explicit} algebraic representatives of every rational Weil
class from a single detected surface.
(2) The forcing input is stratified into three forms --- pair form, carrier form, specified-carrier
form --- and the layers between the forms (a bundle-monoid layer and the $K$-torsion ambiguity) are
identified; a conditional fusion lemma is proved, and the monoid layer is defined as an open stratification
layer (no strictness instance is claimed; on the projective line every pure character is carrier-realizable, so
the monoid question is a higher-dimensional phenomenon).
(3) The End-mass law: on abelian worlds every holomorphic realization of the exact carrier has
$\chi(A,\mathrm{End}\,F)= +m^2\int\alpha^2 >0$ and, on the discriminator, the even-End-cohomology identity
$h^0+\tfrac12 h^2 = h^1 + 4m^2 \ge 4m^2$; consequently \emph{no birth law producing carriers with bounded
even-End-cohomology mass can realize the forcing input at unbounded scales on abelian fourfolds} (the
End-mass no-go; simple realizations at scale $m=1$ satisfy $h^2\ge 6$ on the discriminator). The periods
$\int W_R^2=8$, $\int W_I^2=8$,
$\int W_RW_I=0$, $\Omega\overline\Omega=16\,\mathrm{vol}$ are re-derived exactly by an independent exterior-algebra
computation, and the topological $K$-lattice of every complex torus is proved ch-unimodular, so the
Weil carrier exists at scale $m_0=1$ with fully integral Chern data
$c_1=0,\ c_2=-mW_R,\ c_3=0,\ c_4=4m^2[A]$.
\textbf{What it does not prove:} the Hodge conjecture, its negation, or any reduction of the forcing input to a weaker
statement; the equivalence map of Theorem 2 is recorded as input normalization, not as progress toward the target
(banned-substitute taxonomy M3).}}
\end{center}

\section{Setup: the discriminator and the projector algebra}\label{sec:correspondence}

\subsection{The world}
Throughout, $A$ is a $Q(i)$-Weil abelian fourfold: a polarized abelian fourfold over $\mathbb{C}$ with an action of
$K=\mathbb{Q}(i)$ of Hodge signature $(2,2)$ and a $K$-compatible ample class $\theta\in\mathrm{NS}(A)$, assumed
\emph{nonsplit} (not isogenous to the split CM elliptic power $E^4$) and $\mathrm{NS}(A)_{\mathbb{Q}}=\mathbb{Q}\theta$.
Let $e_1,\dots,e_4,f_1,\dots,f_4\in H^1(A,\mathbb{Z})$ be the standard integral basis with $J e_j=f_j$, $J f_j=-e_j$
($J$ the complex structure), $\Omega=\bigwedge_{j=1}^4(e_j+if_j)\in H^4(A,\mathbb{Z}(i))$, and
$W_R=\mathrm{Re}\,\Omega$, $W_I=\mathrm{Im}\,\Omega$ the two rational real Weil classes. These are rational
$(2,2)$ Hodge classes, primitive against any $K$-compatible polarization.

\subsection{The isogeny and its spectrum}
Let $f=[2+i]\in\mathrm{End}(A)$ (an isogeny). On $H^4(A,\mathbb{Q})$ the operator $f^*$ has the factorization
recorded in the 2026-10-19 Gram record. The polynomial identities of \S\ref{sec:projpoly} are
kernel-verified in the companion certificate file; the spectrum factorization itself is a record-level input
(the eigenvalue census of the 2026-10-19 record), and the idempotence fact (d) is kernel-verified in the
pre-existing \texttt{GSTUniversalCMArithmeticObstruction.lean} given that census:
\[
G_W(t)=t^2+14t+625 \quad(\text{the Weil primary factor, }\deg 2),\qquad
G_{\mathrm{m}}(t)=t^2-30t+625,\qquad (t-25),
\]
where the Weil factor is the minimal polynomial of $(2+i)^4=-7+24i$. The three factors are pairwise coprime over
$\mathbb{Q}$, so $H^4(A,\mathbb{Q})$ is the direct sum of the corresponding primary subspaces
\[
H^4(A,\mathbb{Q}) \;=\; H_W \;\oplus\; H_{\mathrm{m}} \;\oplus\; H_{25},
\]
and the \emph{rational Weil plane} is $W_{\mathbb{Q}}:=H_W=\mathrm{span}_{\mathbb{Q}}\{W_R,W_I\}$ (the 2026-10-19
record: the $16$-coordinate vectors of $\Omega$ have rank $2$ over $\mathbb{Q}$).

\subsection{The projector polynomial and the correspondence}\label{sec:projpoly}
Define
\[
P_W(t)\;=\;-\frac{(t-25)(39t-79)(t^2-30t+625)}{44\,000\,000}.
\]
The following four facts are polynomial identities over $\mathbb{Q}$ ((a)--(c) kernel-verified in the companion
Lean file; (d) kernel-verified in the pre-existing \texttt{GSTUniversalCMArithmeticObstruction.lean}, given the
record-level spectrum census; the last three rest on the key identity below):
\begin{enumerate}[label=(\alph*)]
\item \textbf{Key identity:}\quad $(t-25)(39t-79)(t^2-30t+625)+44\,000\,000=(t^2+14t+625)(39t^2-2770t+72375)$.
\item $P_W\equiv 1 \pmod{G_W}$.
\item $P_W\equiv 0 \pmod{G_{\mathrm{m}}}$ and $P_W\equiv 0\pmod{(t-25)}$.
\item $P_W^2-P_W\equiv 0\pmod{G_WG_{\mathrm{m}}(t-25)}$ (idempotence on the full spectrum), by (b),(c) and the
Chinese remainder theorem for the pairwise coprime factorization.
\end{enumerate}
Expanding the numerator $-(t-25)(39t-79)(t^2-30t+625)=-39t^4+2224t^3-57970t^2+718000t-1234375$ gives the
\textbf{correspondence coefficients}
\begin{equation}\label{eq:coeffs}
P_W(t)\;=\;-\tfrac{39}{44\,000\,000}\,t^4+\tfrac{139}{2\,750\,000}\,t^3-\tfrac{527}{400\,000}\,t^2
+\tfrac{359}{22\,000}\,t-\tfrac{79}{2816}.
\end{equation}

\begin{definition}[The Weil Chow correspondence]\label{def:corr}
Let $\Gamma_{f^j}\in \mathrm{CH}^{4}(A\times A)$ denote the transpose graph of the $j$-th iterate of $f$
(a prime algebraic cycle, $j=0,\dots,4$; $\Gamma_{f^0}$ is the diagonal). Define the rational codimension-$4$
correspondence
\[
\Pi_W^{\mathrm{CH}}\;:=\;\sum_{j=0}^{4}a_j\,\Gamma_{f^j},\qquad
(a_0,\dots,a_4)=\left(-\tfrac{79}{2816},\ \tfrac{359}{22\,000},\ -\tfrac{527}{400\,000},\
\tfrac{139}{2\,750\,000},\ -\tfrac{39}{44\,000\,000}\right),
\]
acting on codimension-$2$ rational Chow classes by
$\Pi_W^{\mathrm{CH}}(Z):=(p_2)_*(\Pi_W^{\mathrm{CH}}\cdot p_1^*[Z])\in \mathrm{CH}^{2}(A)_{\mathbb{Q}}$
and on cohomology by the same formula.
\end{definition}

\begin{lemma}[Cohomology action of the correspondence]\label{lem:action}
On $H^4(A,\mathbb{Q})$, the operator induced by $\Pi_W^{\mathrm{CH}}$ equals $P_W(f^*)$; consequently it is the
identity on $W_{\mathbb{Q}}$, zero on $H_{\mathrm{m}}$ and on the $(t-25)$-primary part $H_{25}$, and it acts on $H^4$ only
(since $S$ is built from degree-$4$ classes, no claim is made about $H^2$ or $H^6$). Moreover $\Pi_W^{\mathrm{CH}}\circ f^{*}=f^{*}\circ\Pi_W^{\mathrm{CH}}$ on Chow and
on cohomology, since both are polynomials in the graph correspondences of $f$.
\end{lemma}

\begin{proof}
The transpose-graph correspondence $\Gamma_{f^j}$ acts on cohomology by $(f^j)^{*}=t^j$ on
$H^{4}$ (standard: $(p_2)_*(\Gamma_{f^j}\cdot p_1^*c)= (f^j)^{*}c$ for all $c\in H^{4}(A,\mathbb{Q})$;
external leaf E1 in \S\ref{sec:scope}). Summing with the coefficients \eqref{eq:coeffs} gives $P_W(f^*)$.
Facts (b),(c) of \S\ref{sec:projpoly} evaluate this on each primary summand: on $H_W$, $f^*$ has minimal polynomial
$G_W$ and $P_W\equiv1$, so $P_W(f^*)=\mathrm{id}$; on $H_{\mathrm{m}}$ and $H_{25}$, $P_W\equiv 0$, so
$P_W(f^*)=0$. The commutation is polynomial composition of commuting graph correspondences
(external leaf E1).
\end{proof}

\section{Theorem W1: the completed zero-or-full frame}\label{sec:W1}

\begin{definition}[The projected Chow image]\label{def:S}
\[
S\;:=\;\Pi_W^{\mathrm{CH}}\big(\mathrm{cl}_A(\mathrm{CH}^{2}(A)_{\mathbb{Q}})\big)\;\subseteq\;H^{4}(A,\mathbb{Q}).
\]
\end{definition}

\begin{theorem}[Geometric zero-or-full dichotomy, completed frame]\label{thm:W1}
Let $A$ be as in \S\ref{sec:correspondence}. Then:
\begin{enumerate}[label=(\roman*)]
\item $S\subseteq W_{\mathbb{Q}}$, and $S$ is stable under $f^{*}$ and $f_{*}$.
\item $S=0$ or $S=W_{\mathbb{Q}}$.
\item $S=0$ implies $W_R$ and $W_I$ are not algebraic, i.e.\ the rational Hodge conjecture fails on $A$
for the classes $W_R,W_I$.
\item $S=W_{\mathbb{Q}}$ if and only if there exists an algebraic surface $Y\subset A$ (equivalently, an integral
codimension-$2$ cycle) with $\Pi_W^{\mathrm{CH}}\mathrm{cl}(Y)\ne 0$; and in that case \emph{every} rational Weil
class is the class of an explicit rational algebraic codimension-$2$ cycle constructed from $Y$ by rational
correspondences.
\item \textbf{(Carrier connection.)} If $S=0$, then for \emph{every} $m\ge 1$ and \emph{every} $r\ge 5$ there is
\emph{no} holomorphic vector bundle $F$ of rank $r$ on $A$ with
$\mathrm{ch}(F)=r+mW_R$ (pure character, all other grades zero); and likewise no pure pair for $mW_R$.
Conversely, if for \emph{some} $m\ge1$, $r\ge5$ a holomorphic $F$ with $\mathrm{ch}(F)=r+mW_R$ exists, then
$S\ne0$, hence $S=W_{\mathbb{Q}}$ and every rational Weil class is algebraic.
\end{enumerate}
\end{theorem}

\begin{proof}
(i) $S$ is generated by classes $\Pi_W^{\mathrm{CH}}\mathrm{cl}(Y)$, each of which lies in the $G_W$-primary
subspace by Lemma \ref{lem:action} (all other primary components are annihilated). Stability under $f^{*}$:
for a cycle $Y$, $f^{*}[Y]\in\mathrm{CH}^{2}(A)_{\mathbb{Q}}$ (pullback of a cycle along the isogeny $f$; external
leaf E2), and $\Pi_W^{\mathrm{CH}}(f^{*}[Y])=f^{*}\Pi_W^{\mathrm{CH}}([Y])$ by Lemma \ref{lem:action}; since $S$ is
by definition generated by the projected classes, this shows $f^{*}S\subseteq S$. The same holds for $f_{*}$ via
$f_{*}\Pi_W^{\mathrm{CH}}=\Pi_W^{\mathrm{CH}}f_{*}$ (graph correspondences commute with pushforward by
finite morphisms, external leaf E2).

(ii) The restriction of $f^{*}$ to $W_{\mathbb{Q}}$ has minimal polynomial $G_W$, which is irreducible over
$\mathbb{Q}$; hence $W_{\mathbb{Q}}$ with its $f^{*}$-action is isomorphic, as a
$\mathbb{Q}[t]$-module, to $\mathbb{Q}[t]/(G_W)\cong(\mathbb{Q}^2,T)$ where $T$ is the companion action
$T(x,y)=(-625y,\,x-14y)$ (kernel-verified as \texttt{cmWeilFourfoldCompanionAction}; both are $2$-dimensional
modules on which $G_W$ is the minimal polynomial, and $\mathbb{Q}[t]/(G_W)$ is the unique such up to isomorphism
because $G_W$ is irreducible of degree $2$). By (i), $S$ is a $T$-invariant $\mathbb{Q}$-submodule of this module,
and the kernel-verified theorem \texttt{cmWeilStableSourceDichotomy} (\texttt{cmWeilInvariantSource\_nonzero\_forces\_full\_plane},
using $\det[z,Tz]=(x-7y)^2+576y^2>0$ for $z\ne0$) gives $S=0$ or $S=W_{\mathbb{Q}}$.

(iii) If $W_R=\mathrm{cl}(Z)$ for a rational cycle $Z$, then $\Pi_W^{\mathrm{CH}}\mathrm{cl}(Z)=W_R\ne0$ by
Lemma \ref{lem:action}, so $S\ne0$; under $S=0$ this is a contradiction. Same for $W_I$, and for every nonzero
element of $W_{\mathbb{Q}}$ (each has nonzero Weil projection: $\Pi_W^{\mathrm{CH}}w=w$ for $w\in W_{\mathbb{Q}}$).

(iv) If $S=W_{\mathbb{Q}}$, pick $z\ne0$ in $S$, say $z=\Pi_W^{\mathrm{CH}}\mathrm{cl}(Y)$ for a rational cycle
$Y$. The kernel-verified one-source reconstruction \texttt{cmWeilOneSourceReconstruction} (with eigenvalue data
$a=-7$, $b=24$) provides explicit rational coefficients $c_0(z,w)$, $c_1(z,w)$ with
$w=c_0(z,w)\,z+c_1(z,w)\,Tz$ for every target $w\in\mathbb{Q}^2$. Under the module identification of (ii), $Tz$
is the class $f^{*}z=\Pi_W^{\mathrm{CH}}\mathrm{cl}(f^{*}Y)$ (again algebraic), so
\[
\mathrm{cl}\big(c_0\,\Pi_W^{\mathrm{CH}}(Y)+c_1\,\Pi_W^{\mathrm{CH}}(f^{*}Y)\big)=w,
\]
an explicit rational algebraic representative of $w$. The converse is immediate from the definition of $S$.

(v) Suppose $F$ holomorphic of rank $r\ge5$ with $\mathrm{ch}(F)=r+mW_R$ exactly. By GAGA (external leaf E3) $F$ is
algebraic; fix $k\gg0$ so that $F\otimes O(k)$ is globally generated, and let $D\subset A$ be the
maximal-degeneracy scheme of a general map $O^{r-1}\to F\otimes O(k)$ from the \emph{trivial} source
(the cycle-1 convention; its section space $H^0(F\otimes O(k))^{\oplus(r-1)}$ grows with $k$), which is of pure
codimension $2$ with class $\mathrm{c}_2^{\mathrm{CH}}(F\otimes O(k))$ by Thom--Porteous (external leaf E4, at maximal
degeneracy, expected codimension $2$). The two-center descent law in single-center form (proved for all $p$ in the
2026-10-24 lattice completion; $p=2$ here) gives, with the splitting-principle identity
$\mathrm{c}_2^{\mathrm{CH}}(F\otimes O(k))=\mathrm{c}_2^{\mathrm{CH}}(F)+\binom{r}{1}k\theta\,\mathrm{c}_1+\binom{r}{2}k^2\theta^2$
and the carrier equations $\mathrm{c}_1=0$, $\mathrm{c}_2(F)=-mW_R$ (Newton conductor lemma, same record):
\[
\big[D\big]\;=\;-mW_R+\tbinom{r}{2}k^{2}\theta^{2}\quad\text{in }\mathrm{CH}^{2}(A)_{\mathbb{Q}},
\]
hence $Z:=\tfrac1m\big(\binom{r}{2}k^{2}\theta^{2}-[D]\big)$ is an explicit rational algebraic cycle with
$\mathrm{cl}(Z)=W_R$, so $S\ne0$ --- contradiction with $S=0$. The identical argument with the pair form
(Theorem C, pair version) rules out pure pairs. The converse is (iii) read contrapositively: a holomorphic carrier
gives an algebraic $W_R$, hence $S\ne0$, hence (ii) $S=W_{\mathbb{Q}}$, and (iv) applies.
\end{proof}

\begin{remark}[What the frame does and does not decide]
Theorem \ref{thm:W1} converts the discriminator into a single well-posed binary geometric question:
does one algebraic surface with nonzero Weil projection exist on $A$? Both outcomes are now closed, kernel-anchored
statements about their consequences. It does not decide which outcome holds on any actual nonsplit $A$; producing
$S=0$ on an instance is the campaign's negative certificate (PROMPT \S2), and the 2026-10-18/19 records'
divisor-channel no-go ($\Pi_W(\theta^p)=0$) verifies $S=0$ only on the divisor-generated subalgebra, not on all
surfaces. The external manuscript claim (\texttt{p3sw:weil-plane}, unauthenticated in-repo) would, if ever verified,
force the $S=W_{\mathbb{Q}}$ outcome; the campaign treats it as an unverified external leaf.
\end{remark}

\section{Theorem W2: the three forms of the forcing input}\label{sec:W2}

The forcing input for a class $\alpha\in H^{2p}(X,\mathbb{Q})$ of type $(p,p)$ (with $m\equiv0 \pmod{m_0}$,
$r\ge\max(p+1,d+1)$) has appeared in the campaign in three forms. This section fixes their exact
relations. \emph{Scope note: this is a normalization map of known equivalences and open layers --- by the
banned-substitute taxonomy M3 it is not progress toward the target; its value is that it identifies which layer
each future mechanism addresses.}

\begin{definition}[The three input forms]
For a pure character $\mathrm{ch}=r+m\alpha$ (grade $0$: $r$; grade $p$: $m\alpha$; all other grades zero):
\begin{enumerate}[label=(\arabic*)]
\item \textbf{Pair form (at some scale):} there exist an integer $m'\ge1$ and $x\in K^0_{\mathrm{alg}}(X)$ with
$\mathrm{ch}(x)=m'\alpha$ (grade $p$ only, $\mathrm{rank}(x)=0$). At a pinned scale $m$ we say \emph{pair form at scale $m$}.
\item \textbf{Carrier form:} there exists a holomorphic vector bundle $F$ of rank $r\ge\max(p+1,d+1)$ with
$\mathrm{ch}(F)=r+m\alpha$.
\item \textbf{Specified-carrier form:} the specific topological bundle $E_\alpha$ of Theorem A admits an integrable
$\bar\partial$-operator.
\end{enumerate}
\end{definition}

\begin{lemma}[Trivial directions]\label{lem:trivial}
$(3)\Rightarrow(2)\Rightarrow(1)$. Moreover \emph{pair form at some positive integer scale} $\Leftrightarrow$
$\alpha$ is algebraic: ($\Rightarrow$) apply the two-center descent law in pair form to a bundle presentation of $x$
(equal ranks and equal low Chern classes follow from purity by the Newton conductor lemma), which returns a
cycle of class $\alpha$ after dividing by $m'$; ($\Leftarrow$) from a cycle with class $\alpha$, the Theorem
B2(ii) construction returns an algebraic pure K-class at \emph{some} integer scale (the pinned-scale statement
requires an additional saturation step which is neither claimed nor needed here; cf.\ the exactness note in the
scope section).
\end{lemma}

\begin{lemma}[Conditional fusion: presentable pairs give carriers]\label{lem:fusion}
Let $x\in K^0_{\mathrm{alg}}(X)$ with $\mathrm{ch}(x)=m\alpha$ pure, and suppose $x$ admits a presentation
$x=[E_+]-[E_-]$ by holomorphic bundles of equal rank $s$ together with a surjection
$(E_+\oplus O^{r})\otimes O(k)\twoheadrightarrow E_-\otimes O(k)$ for some $k$ (a \emph{one-step presentable} pair).
Then $F:=\ker(A)\otimes O(-k)$ is a holomorphic bundle of rank $r$ with
\[
\mathrm{ch}(F)\;=\;\mathrm{ch}(E_+)+r-\mathrm{ch}(E_-)\;=\;r+m\alpha .
\]
\end{lemma}

\begin{proof}
The kernel of a surjective map of locally free sheaves is locally free (locally split), of rank $s+r-s=r$.
Chern characters are additive on short exact sequences and satisfy
$\mathrm{ch}(G\otimes L)=\mathrm{ch}(G)\mathrm{ch}(L)$ for line bundles $L$ (external leaf E5), so
\[
\mathrm{ch}(F)=\mathrm{ch}(\ker A)\mathrm{ch}(O(-k))
=\big(\mathrm{ch}(E_+)\mathrm{ch}(O(k))+r\,\mathrm{ch}(O(k))-\mathrm{ch}(E_-)\mathrm{ch}(O(k))\big)\mathrm{ch}(O(-k))
=(\mathrm{ch}(E_+)+r-\mathrm{ch}(E_-)),
\]
which is the pure character by the presentation. Rank: $s+r-s=r$.
\end{proof}

\begin{remark}[Why the fusion is only conditional]
A surjection $(E_+\oplus O^{r})\otimes L\twoheadrightarrow E_-\otimes L$ is equivalent (after tensoring by $L^\vee$)
to a surjection $E_+\oplus O^{r}\twoheadrightarrow E_-$, and the space of such maps is the \emph{fixed}
space $H^0((E_+\oplus O^r)^\vee\otimes E_-)$: amplification does not enlarge it. In general a given pair carries no
such surjection, so no unconditional fusion follows from this argument. The computation after this remark shows
that on the projective line the monoid layer is in fact empty --- the layer question is strictly
higher-dimensional.
\end{remark}

\begin{proposition}[Monoid computation on the projective line]\label{prop:monoid}
On $X=\mathbb{P}^1$ with $L=[O(1)]$ and $K^0(\mathbb{P}^1)=\mathbb{Z}[L]/((L-1)^2)$:
\begin{enumerate}[label=(\alph*)]
\item $[O(k)]=[O]+k\,([O(1)]-[O])$ for every $k\in\mathbb{Z}$ (since $(L-1)^2=0$, $L^k=1+k(L-1)$).
\item Consequently every rank-$r$ class $r+q\,(L-1)$ ($q\in\mathbb{Z}$, $r\ge1$) is the class of the bundle
$O(q)\oplus O^{\oplus(r-1)}$: the projective line is \emph{monoid-full} (every stable-rank algebraic K-class is a
single bundle class, and in particular every pure character $r+m\theta$ is carrier-realizable for every
$m\in\mathbb{Z}$, by $q=m$). The monoid layer is therefore a strictly higher-dimensional phenomenon; no
strictness instance is claimed on any world in this document.
\end{enumerate}
\end{proposition}

\begin{proof}
(a) $L^k=(1+(L-1))^k=1+k(L-1)$ because $(L-1)^2=0$. (b) $[O(q)\oplus O^{r-1}]=L^{q}+(r-1)=1+q(L-1)+r-1=
r+q(L-1)$, and every rank-$r$ class has the form $r+q(L-1)$ since $(L-1)$ spans the reduced group. For the pure
character: $\mathrm{ch}(O(m)\oplus O^{r-1})=e^{m\theta}+r-1=r+m\theta$ because $\theta^2=0$ on $\mathbb{P}^1$.
\end{proof}

\begin{theorem}[Input stratification]\label{thm:W2}
Fix $(X,p,\alpha)$ with $m\equiv0\pmod{m_0}$, $r\ge\max(p+1,d+1)$.
\begin{enumerate}[label=(\roman*)]
\item Pair form at some positive integer scale $\Leftrightarrow$ $\alpha$ algebraic (Lemma \ref{lem:trivial}).
\item Carrier form $\Rightarrow$ pair form; specified-carrier form $\Rightarrow$ carrier form. The two layers between
the forms are well-defined:
\begin{enumerate}[label=(\roman*)]
\item the \textbf{monoid layer}: write $\mathcal{M}(X)$ for the inner set $\{$pure characters realized by single
holomorphic bundles$\}$, sitting inside $\{$pure characters realized by algebraic $K$-classes$\}$; the containment
is an \emph{open} layer question: it is an equality on monoid-full worlds (Proposition \ref{prop:monoid}: the
projective line is monoid-full), and no world on which it is strict is exhibited here;
\item the \textbf{torsion gap}: two specified carriers with the same character differ by an element of
$\ker(\mathrm{ch})=K^0(X)_{\mathrm{tors}}$ (finitely generated $K^0$, $\mathrm{ch}\otimes\mathbb{Q}$ injective;
external leaf E7), so specified-carrier form differs from carrier form by Theorem A's torsion ambiguity
(a qualitative gap; the quotient $T(K^0)/T(K^0)\cap K^0_{\mathrm{alg}})$ measures one part of it, up to the
monoid question).
\end{enumerate}
\item On the discriminator $A$: the pair form at scale $m$ implies $S\ne0$; the carrier form implies $S\ne0$
(Theorem \ref{thm:W1}(v)); and under $S=W_\mathbb{Q}$ the pair form holds at an explicit scale (the one-source
coefficients clear a finite denominator). The residual monoid layer
\[
\mathcal{M}(A,W):=\{\text{carrier-realizable pure Weil characters}\}
\]
is a well-defined open layer of the forcing site on the discriminator, separating the pair outcome from the
carrier outcome (emptiness is not claimed).
\end{enumerate}
\end{theorem}

\begin{proof}
(i) is Lemma \ref{lem:trivial}. (ii): the forward implications hold because a holomorphic bundle is algebraic
(GAGA) and its class is an algebraic K-class, and because the specified carrier is by definition a bundle. The
monoid layer: Proposition \ref{prop:monoid} proves the projective line \emph{monoid-full}; whether any world has a
nonempty layer is left open --- the layer is defined, not instantiated. The torsion gap:
$K^0(X)$ is finitely generated (finite CW type) and $\mathrm{ch}\otimes\mathbb{Q}$ is injective (Atiyah--Hirzebruch,
external leaf E7), so $\ker(\mathrm{ch})$ is exactly the torsion subgroup; Theorem A constructs the carrier
uniquely up to this subgroup, which measures (qualitatively, together with the monoid question) the difference
between the specified-carrier and carrier forms. (iii): Theorem \ref{thm:W1}(v) and its proof. For the explicit
scale under $S=W_\mathbb{Q}$: write the detected representatives as rational combinations $\sum_i q_i[Z_i]$ with
$\mathrm{cl}$-class $n_0W$ (the one-source reconstruction; $n_0$ the lcm of its denominators). On the abelian $A$
with $\mathrm{td}(A)=1$, GRR (external leaf E8) gives $\mathrm{ch}([O_{Z_i}])=[Z_i]+\iota_*\mathrm{td}_1(Z_i)+\iota_*\mathrm{td}_2(Z_i)$:
the grade-$3$ and grade-$4$ junk are pushforwards of algebraic classes, hence algebraic; correct them by
subtracting $\sum_j s_j[O_{C_j}]$ for algebraic curves $C_j$ with the right codimension-$3$ classes and then a
rational multiple of $[O_{\mathrm{pt}}]$ (whose character is the pure top class). The result is a rational algebraic K-class $x\in K^0_{\mathrm{alg}}(A)\otimes\mathbb{Q}$ with $\mathrm{ch}(x)=W$ pure;
after clearing denominators $m:=$ the lcm of the denominators of the $q_i$, $s_j$ and the point-correction, the
integral class $m\cdot x\in K^0_{\mathrm{alg}}(A)$ has $\mathrm{ch}=mW$ pure --- pair form at an explicit scale, as claimed.
\end{proof}

\begin{remark}[Effect on the cycle-1 audit flag]
The 2026-10-24 audit (finding F4) flagged the registry's forcing-site biconditional as overstated because
``HC gives algebraic realization of the K-class only up to K-torsion''. Theorem \ref{thm:W2} resolves the flag
into its true shape: the pair form is equivalent to algebraicity (no torsion loss); the torsion gap separates
carrier from specified carrier; and a further, previously untracked monoid gap separates pair from carrier. The
descent law consumes the \emph{character} only, so Theorem A's torsion ambiguity was already harmless for the
cycle-1 construction; the monoid layer is the substantive new boundary.
\end{remark}

\section{Theorem W3: the End-mass law}\label{sec:endsignature}

\begin{lemma}[Endomorphism character of the exact carrier]\label{lem:endchar}
Let $E$ be any complex vector bundle (topological or holomorphic) with
$\mathrm{ch}(E)=r+m\alpha$ pure (grade $p$ only). Then
\[
\mathrm{ch}(E^\vee)=r+(-1)^p m\alpha \quad\text{and}\quad
\mathrm{ch}(\mathrm{End}\,E)=\mathrm{ch}(E)\,\mathrm{ch}(E^\vee)
= r^{2}+\bigl(1+(-1)^p\bigr)rm\alpha+(-1)^p m^{2}\alpha^{2},
\]
so the endomorphism character has three grades in general: $r^2$ in grade $0$, $2rm\alpha$ in grade $p$
(\emph{only for $p$ even}; for $p$ odd this grade cancels), and $(-1)^p m^2\alpha^2$ in grade $2p$. At $p=2$:
$\mathrm{ch}(\mathrm{End}\,E)=r^{2}+2rm\alpha+m^{2}\alpha^{2}$.
\end{lemma}

\begin{proof}
$\mathrm{ch}(E^\vee)=\sum_j(-1)^{j}\mathrm{ch}_j(E)$ (dual roots $e^{-x_i}$; external leaf E5): the $p$-th grade
picks up the sign $(-1)^p$ --- at even $p$ the dual character COINCIDES with the original, it does not negate it.
$\mathrm{ch}(E\otimes G)=\mathrm{ch}(E)\mathrm{ch}(G)$ (external leaf E5), with $\mathrm{End}\,E=E\otimes E^\vee$.
Multiply and collect grades: grade $0$: $r\cdot r$; grade $p$: $r(-1)^pm\alpha+m\alpha\cdot r$; grade $2p$:
$(-1)^pm^2\alpha^2$.
\end{proof}

\begin{remark}[The parity trap this lemma repairs]
A previous version of this section used $\mathrm{ch}(E^\vee)=r-m\alpha$ at $p=2$, i.e.\ applied the sign $-1$
globally instead of $(-1)^p$ per grade; that is the odd-$p$ identity, false at even $p$, and it flipped the sign of
$\chi(\mathrm{End}\,F)$ below. The corrected computation is the point of Theorem \ref{thm:W3}. The same note applies
to the cycle-1 registry's queued ``pure grade-$2p$ HR-positive End handle'' --- at even $p$ that handle does not
exist in the stated form.
\end{remark}

\begin{theorem}[End-mass law]\label{thm:W3}
Let $X$ be a smooth projective fourfold, $\alpha$ a nonzero real primitive $(2,2)$ class, $p=2$, and suppose
$F$ is \emph{any} holomorphic realization of the exact carrier character
$\mathrm{ch}(F)=r+m\alpha$ ($r\ge5$, $m\ge1$). Then, using Lemma \ref{lem:endchar} at $p=2$
($\mathrm{ch}(\mathrm{End}\,F)=r^2+2rm\alpha+m^2\alpha^2$):
\begin{enumerate}[label=(\roman*)]
\item (Riemann--Roch, external leaf E8; $\mathrm{td}(A)=1$ on abelian $A$):
on an abelian fourfold $A$,
\[
\chi(A,\mathrm{End}\,F)\;=\;\int_A \mathrm{ch}_4(\mathrm{End}\,F)\;\;=\;\;+\,m^2\int_A\alpha^{2}\;>\;0,
\]
(the $r^2$-term has degree $0$ and pairs only with $\mathrm{td}_4$, whose integral is $\chi(O_A)=0$;
the $2rm\alpha$-term has degree $4$ and pairs only with $\mathrm{td}_2=0$; the degree-$8$ part is $m^2\alpha^2$).
(ii$_{\mathrm{abelian}}$) On abelian $A$ (Serre duality, external leaf E9; $K_A$ trivial):
$h^3=h^1$, $h^4=h^0$, so $\chi=2h^0+h^2-2h^1$, and substituting (i):
\[
h^0(\mathrm{End}\,F)+\tfrac12 h^2(\mathrm{End}\,F)
\;=\;h^1(\mathrm{End}\,F)+\tfrac{m^2}{2}\int_A\alpha^2
\;\ge\;\tfrac{m^2}{2}\int_A\alpha^2 ,
\]
using $h^1\ge0$. Equivalently $h^1=h^0+\tfrac12h^2-\tfrac{m^2}{2}\int\alpha^2$.
(i$_{\mathrm{gen}}$) On a general smooth projective fourfold the formula reads
$\chi(\mathrm{End}\,F)=r^2\chi(O_X)+2rm\int_X\alpha\,\mathrm{td}_2(X)+m^2\int_X\alpha^2$.
\end{enumerate}
\end{theorem}

\begin{proof}
(i) Riemann--Roch for the bundle $\mathrm{End}\,F$ on the fourfold:
$\chi=\int\mathrm{ch}(\mathrm{End}F)\mathrm{td}(X)$; by Lemma \ref{lem:endchar} the character is
$r^2+2rm\alpha+m^2\alpha^2$, so on an abelian fourfold $\mathrm{td}=1$ (all rational Chern classes of the
tangent bundle vanish) and the degree-$8$ part of the integrand is exactly $m^2\alpha^2$; the displayed pairing
analysis gives $\chi=+m^2\int\alpha^2>0$ (the period $\int\alpha^2>0$ is the verified Hodge--Riemann-positivity
input; on the discriminator it is the certified period $\int W_R^2=8$).
(ii$_{\mathrm{abelian}}$) Serre duality on the abelian variety: $h^i(\mathrm{End}F)=h^{4-i}(\mathrm{End}F\otimes K_A)
=h^{4-i}(\mathrm{End}F)$, so $h^4=h^0$ and $h^3=h^1$ (using $\mathrm{End}^\vee\cong\mathrm{End}$). Expanding
$\chi=h^0-h^1+h^2-h^3+h^4=2h^0+h^2-2h^1$ and substituting (i) gives
$2h^0+h^2-2h^1=m^2\int\alpha^2$, the displayed identity; $h^1\ge0$ gives the bound.
(i$_{\mathrm{gen}}$): on a general fourfold the degree-$8$ part of
$(r^2+2rm\alpha+m^2\alpha^2)\mathrm{td}(X)$ is $r^2\mathrm{td}_4+2rm\alpha\,\mathrm{td}_2+m^2\alpha^2$, and
$\int\mathrm{td}_4=\chi(O_X)$ (external leaf E8, Hirzebruch--Riemann--Roch).
\end{proof}

\begin{corollary}[End-mass no-go for abelian worlds]\label{cor:rigidity}
On any abelian fourfold $A$ and any nonzero real primitive $(2,2)$ class $\alpha$ with $\int_A\alpha^2>0$
(e.g.\ the Weil classes: $\int W_R^2=8$ by \S\ref{sec:periods}): every holomorphic realization of the exact
carrier at scale $m$ satisfies
\[
h^0(A,\mathrm{End}\,F)+\tfrac12 h^2(A,\mathrm{End}\,F)\;\ge\;\tfrac{m^2}{2}\int_A\alpha^2,
\]
which equals $4m^2$ on the discriminator; for \emph{simple} realizations ($h^0=1$) this forces
$h^2(\mathrm{End}\,F)\ge 8m^2-2$ (in particular $h^2\ge6$ already at the minimal scale $m=1$ on the
discriminator). Consequently \emph{any} birth law that produces carriers whose even-End-cohomology mass
$h^0+h^2/2$ stays bounded independently of $m$ cannot realize the forcing input at all scales on abelian
fourfolds. No claim is made about rigidity or isolatedness: the corrected identity shows $h^1$ can be $0$ with
the mass carried by $h^2$.
\end{corollary}

\begin{remark}[Relation to the polystable no-go]
The 2026-10-18 record (universal determinantal packet interpolation) proved that a slope-polystable holomorphic
bundle with $c_1=0$ and $c_2=\pm N\alpha$ violates the Chern--Weil zero-energy identity whenever
$\int\alpha^2>0$ and $\alpha\theta^{d-2}=0$: the polystable channel is dead for primitive carriers. Theorem
\ref{thm:W3} is complementary: it constrains whatever \emph{survives} --- any realization must be unstable and
must carry even-End-cohomology mass at least $\tfrac{m^2}{2}\int\alpha^2$ ($4m^2$ on the discriminator, so at
the minimal scale $m=1$ every simple realization has $h^2\ge6$). Both constraints are blind to the
Hodge/non-Hodge distinction: they apply to every holomorphic bundle
with the carrier's Chern data, and thus do not decide the dichotomy of Theorem \ref{thm:W1}; they prune mechanism
families.
\end{remark}

\section{Exact periods and the minimal scale on the discriminator}\label{sec:periods}

\subsection{The exterior computation}
Work in the exterior algebra on $e_1,\dots,e_4,f_1,\dots,f_4$ (indices $1..8$), canonical order
$e_1e_2e_3e_4f_1f_2f_3f_4$, volume $\mathrm{vol}=e_1e_2e_3e_4f_1f_2f_3f_4$. Expanding
$\Omega=\bigwedge_{j=1}^4(e_j+if_j)$ and taking real parts (even $|S|$; $i^{|S|}=(-1)^{|S|/2}$):
\[
W_R=\;\underset{S=\emptyset}{e_{1234}}
\;-\underset{S=\{3,4\}}{e_{12}f_{34}}
\;+\underset{S=\{2,4\}}{e_{13}f_{24}}
\;-\underset{S=\{2,3\}}{e_{14}f_{23}}
\;-\underset{S=\{1,4\}}{e_{23}f_{14}}
\;+\underset{S=\{1,3\}}{e_{24}f_{13}}
\;-\underset{S=\{1,2\}}{e_{34}f_{12}}
\;+\underset{S=\{1,2,3,4\}}{f_{1234}},
\]
an explicit $8$-term $\pm1$ combination (the signs are the product of $(-1)^{|S|/2}$ with the reordering sign to
canonical form; verified term-by-term in the companion computation and in the Lean certificate as data). Each
pair of complementary even subsets contributes to $W_R^2$ exactly twice (ordered pairs), all non-complementary
products vanish by repeated indices, and the eight ordered contributions carry signs whose total is $+8$:
\begin{equation}\label{eq:WR2}
W_R^2\;=\;8\,\mathrm{vol}\qquad\text{(as a class: no other nonzero component),}
\end{equation}
and symmetrically $W_I^2=8\,\mathrm{vol}$, while the mixed products cancel to $W_RW_I=0$ and
$\Omega\overline\Omega=W_R^2+W_I^2=16\,\mathrm{vol}$. This reproduces the 2026-10-19 record's Wolfram computation
by an independent method.

\subsection{Unimodularity of the torus $K$-lattice and the minimal scale}

\begin{lemma}[ch-unimodularity on tori]\label{lem:unimod}
Let $T$ be a complex torus of dimension $n$. Write $T=(T_2)^{n}$ with $T_2$ the polarized real $2$-torus. The
K\"unneth basis of $K^0(T)$ --- tensor products over the $T_2$-factors of the classes $1,\eta=[O]-[L]$ (even
$K^0$-factors) and of the two odd generators of $K^1(T_2)$ (pairs of even numbers of odd factors) --- maps under
$\mathrm{ch}$ to $\pm$ the standard integral monomial basis of $H^{\mathrm{even}}(T,\mathbb{Z})$:
on each $T_2$-factor $\mathrm{ch}(\eta)=-x$ (the top class; $e^{x}=1+x$ since $x^2=0$ on $T_2$) and
$\mathrm{ch}(\text{odd generator})=\pm$ (the corresponding degree-$1$ class). Hence
$\mathrm{ch}\colon K^0(T)\to H^{\mathrm{even}}(T,\mathbb{Z})$ is surjective and carries the lattice onto the lattice.
\end{lemma}

\begin{proof}
K\"unneth for complex $K$-theory over the product decomposition (external leaf E10) identifies $K^0(T)$ with the
direct sum over parity-compatible tensor products; on each factor the stated Chern character is the exponential of
the degree-$1$ classes, with $x^2=0$ on the two-torus making $\mathrm{ch}(\eta)=1-e^{x}=-x$ exactly, and the odd
Chern character on $K^1$ of the circle factor is the fundamental class. Tensor multiplicativity of $\mathrm{ch}$
completes the identification; signs are permutation signs only.
\end{proof}

\begin{corollary}[Minimal scale on abelian worlds]\label{cor:scale}
On an abelian variety $A$ the integrality scale $m_0$ of Theorem A for a class $\alpha$ equals the denominator
scale of $\alpha$ itself in $H^{2p}(A,\mathbb{Z})$: if $\alpha$ is integral then $m_0=1$ and the carrier exists at
$m=1$. On the Weil discriminator: $W_R$ is integral, $\mathrm{td}(A)=1$, the additive congruence scale of
Theorem B2(i) is $1$ ($c(y)=m\int\alpha\,\mathrm{ch}(y)\mathrm{td}(A)$ is integral for integral $m\alpha$), and the
grade-$4$ Newton class $c_4(F)=\tfrac{m^2}{2}W_R^2=4m^2[A]$ is integral for every $m$. Hence the exact Weil carrier
exists topologically at every scale $m\ge1$ with integral Chern data
\[
c_1=0,\qquad c_2=-mW_R,\qquad c_3=0,\qquad c_4=4m^2[A],\qquad r\ge\max(3,5)=5.
\]
\end{corollary}

\begin{remark}[Consequences for the frame]
At scale $m=1$ the discriminator's forcing input in carrier form is the holomorphic-bundle question for the
character $r+W_R$ with $c_2=-W_R$, and Theorem \ref{thm:W3} gives every realization the signature
the End-mass $h^0+\tfrac12h^2=h^1+4\ge4$ (with $h^2\ge6$ for simple realizations). The two statements together are the campaign's first pair of sharp, world-specific
constraints on the forcing input on its fixed discriminating test.
\end{remark}

\section{Scope, external leaves, and CI}\label{sec:scope}

\textbf{Proved here with full proofs:} Lemma \ref{lem:action}; Theorem \ref{thm:W1} (all five parts, on the stated
external leaves); Lemma \ref{lem:trivial}; Lemma \ref{lem:fusion}; Proposition \ref{prop:monoid}; Theorem
\ref{thm:W2}; Lemma \ref{lem:endchar}; Theorem \ref{thm:W3}; Corollary \ref{cor:rigidity}; the exterior computation
\eqref{eq:WR2} (term-by-term in the companion computation; certified in Lean as data plus verified sign-count
arithmetic); Lemma \ref{lem:unimod}; Corollary \ref{cor:scale}.

\textbf{Kernel-verified companions (GitHub Actions, Lean, sorryAx-audited, \#print axioms):} the projector
polynomial identities (a)--(c) of \S\ref{sec:projpoly} with the explicit coefficients \eqref{eq:coeffs} (the
idempotence (d) is kernel-verified in the pre-existing \texttt{GSTUniversalCMArithmeticObstruction.lean}, given the
record-level spectrum census); the exterior sign-count certificate $\sum(\text{8 ordered contributions})=8$
feeding the End-mass law's discriminator constant; the arithmetic chain from the periods to
$\chi=+8m^2$ and the End-mass identity $h^0+\tfrac12h^2=h^1+4m^2\ge4m^2$.
Data-versus-proved labeling inside the Lean docstring follows the audited convention of
\texttt{GSTGr25PurePairLatticeCertificate.lean}.

\textbf{External leaves (standard, hypotheses stated where used):}
E1: cycle-class compatibility of correspondences (transpose-graph action $=(f^j)^*$; functoriality and commutation
of graph correspondences) --- standard intersection theory (Fulton, \S16--17 conventions).
E2: isogeny pullback/pushforward on rational Chow groups (finite morphism Gysin maps, $f^*f_*=\deg$) --- standard.
E3: GAGA (holomorphic $=$ algebraic on projective $X$; coherent sheaves and Chern classes).
E4: Thom--Porteous at maximal degeneracy with the expected-codimension genericity (used through the already-proved
two-center descent law of 2026-10-24).
E5: $\mathrm{ch}$ multiplicativity for tensor products and line-bundle twists; $\mathrm{ch}(E^\vee)$; additivity on
exact sequences (all used within the already-audited descent-law infrastructure).
E6: the projective-line computation of \S\ref{sec:W2} uses only $(L-1)^2=0$ in $K^0(\mathbb{P}^1)=\mathbb{Z}[L]/((L-1)^2)$
(no Birkhoff--Grothendieck needed after the repair).
E7: finite generation of $K^0$ for finite CW type; Atiyah--Hirzebruch rational injectivity/surjectivity of
$\mathrm{ch}$ (used in Theorem A already).
E8: Hirzebruch--Riemann--Roch for $\mathrm{End}F$; $\int\mathrm{td}_4=\chi(O_X)$ (and $\int\mathrm{td}_2$ where cited).
E9: Serre duality ($h^i(F)=h^{d-i}(F^\vee\otimes K)$); $\mathrm{End}^\vee\cong\mathrm{End}$; $K_A$ trivial.
E10: K\"unneth formula in complex $K$-theory and multiplicativity of $\mathrm{ch}$; odd Chern character on circle
factors.
E11: Hodge--Riemann positivity of the primitive middle square ($\int\alpha^2>0$ for nonzero real primitive
$(2,2)$ classes on a fourfold); used only in Theorem \ref{thm:W3}(i)'s strict positivity (the corollary takes
$\int\alpha^2>0$ as a hypothesis).

\textbf{Explicitly not claimed:} any construction of a holomorphic structure on any Weil carrier; any decision of
$S=0$ versus $S=W_\mathbb{Q}$ on any actual nonsplit instance; any improvement of the cycle-1 conditional landing;
that the monoid layer $\mathcal{M}(A,W)$ is empty or nonempty. The Hodge conjecture itself remains the campaign's
target; the forcing input remains equivalent to it in pair form; this document's contributions are the completed
dichotomy frame, the stratification of the input, the End-mass law and its no-go corollary, and the exact
discriminator constants.

\vspace{2mm}
\noindent\rule{\textwidth}{0.4pt}
\vspace{1mm}

\noindent\textbf{Pre-commit terminology gate (ledger 076):} this document uses the permanent campaign terminology
(``exact carrier'', ``exactness note''); the banned-word list of rule 50 has been grepped and cleared.

\end{document}
